\documentclass[10pt,a4paper]{amsart}
\usepackage{amsmath,amssymb,amsthm,graphicx}
% For commutative diagrams we use Xy-pic
\usepackage[all]{xy}

% We use 2cell for 2-commutative diagrams.
\xyoption{2cell}
% Theorem environments.
%
\theoremstyle{plain}
\newtheorem{theorem}[subsection]{Theorem}
\newtheorem{proposition}[subsection]{Proposition}
\newtheorem{lemma}[subsection]{Lemma}

\theoremstyle{definition}
\newtheorem{definition}[subsection]{Definition}
\newtheorem{example}[subsection]{Example}
\newtheorem{exercise}[subsection]{Exercise}
\newtheorem{situation}[subsection]{Situation}

\theoremstyle{remark}
\newtheorem{remark}[subsection]{Remark}
\newtheorem{remarks}[subsection]{Remarks}

\numberwithin{equation}{subsection}

% Macros
%
\def\lim{\mathop{\mathrm{lim}}\nolimits}
\def\colim{\mathop{\mathrm{colim}}\nolimits}
\def\Spec{\mathop{\mathrm{Spec}}}
\def\Hom{\mathop{\mathrm{Hom}}\nolimits}
\def\Ext{\mathop{\mathrm{Ext}}\nolimits}
\def\SheafHom{\mathop{\mathcal{H}\!\mathit{om}}\nolimits}
\def\SheafExt{\mathop{\mathcal{E}\!\mathit{xt}}\nolimits}
\def\Sch{\mathit{Sch}}
\def\Mor{\mathop{\mathrm{Mor}}\nolimits}
\def\Ob{\mathop{\mathrm{Ob}}\nolimits}
\def\Sh{\mathop{\mathit{Sh}}\nolimits}
\def\NL{\mathop{N\!L}\nolimits}
\def\CH{\mathop{\mathrm{CH}}\nolimits}
\def\proetale{{pro\text{-}\acute{e}tale}}
\def\etale{{\acute{e}tale}}
\def\QCoh{\mathit{QCoh}}
\def\Ker{\mathop{\mathrm{Ker}}}
\def\Im{\mathop{\mathrm{Im}}}
\def\Coker{\mathop{\mathrm{Coker}}}
\def\Coim{\mathop{\mathrm{Coim}}}

% Boxtimes
%
\DeclareMathSymbol{\boxtimes}{\mathbin}{AMSa}{"02}

%
% Macros for moduli stacks/spaces
%
\def\QCohstack{\mathcal{QC}\!\mathit{oh}}
\def\Cohstack{\mathcal{C}\!\mathit{oh}}
\def\Spacesstack{\mathcal{S}\!\mathit{paces}}
\def\Quotfunctor{\mathrm{Quot}}
\def\Hilbfunctor{\mathrm{Hilb}}
\def\Curvesstack{\mathcal{C}\!\mathit{urves}}
\def\Polarizedstack{\mathcal{P}\!\mathit{olarized}}
\def\Complexesstack{\mathcal{C}\!\mathit{omplexes}}
% \Pic is the operator that assigns to X its picard group, usage \Pic(X)
% \Picardstack_{X/B} denotes the Picard stack of X over B
% \Picardfunctor_{X/B} denotes the Picard functor of X over B
\def\Pic{\mathop{\mathrm{Pic}}\nolimits}
\def\Picardstack{\mathcal{P}\!\mathit{ic}}
\def\Picardfunctor{\mathrm{Pic}}
\def\Deformationcategory{\mathcal{D}\!\mathit{ef}}
\newenvironment{reference}{}{}
\newenvironment{slogan}{}{}
\newenvironment{history}{}{}
\usepackage{cleveref}
\crefname{originalsection}{original section}{original sections}
\Crefname{originalsection}{Original section}{Original sections}
\makeatletter
\newlabel{situation-regular-model}{{0C61}{source}{}{sourcepointer-1}{}}
\newlabel{situation-regular-model@cref}{{[originalsection][0][]0C61}{[1][0][]source}{}{}{}}
\makeatother
\crefname{originallemma}{original lemma}{original lemmas}
\Crefname{originallemma}{Original lemma}{Original lemmas}
\makeatletter
\newlabel{lemma-add-component}{{0CA2}{source}{}{sourcepointer-2}{}}
\newlabel{lemma-add-component@cref}{{[originallemma][0][]0CA2}{[1][0][]source}{}{}{}}
\makeatother
\crefname{originallemma}{original lemma}{original lemmas}
\Crefname{originallemma}{Original lemma}{Original lemmas}
\makeatletter
\newlabel{lemma-properties-form}{{0C66}{source}{}{sourcepointer-3}{}}
\newlabel{lemma-properties-form@cref}{{[originallemma][0][]0C66}{[1][0][]source}{}{}{}}
\makeatother
\crefname{originallemma}{original lemma}{original lemmas}
\Crefname{originallemma}{Original lemma}{Original lemmas}
\makeatletter
\newlabel{perfect-lemma-chi-locally-constant-geometric}{{0B9T}{source}{}{sourcepointer-4}{}}
\newlabel{perfect-lemma-chi-locally-constant-geometric@cref}{{[originallemma][0][]0B9T}{[1][0][]source}{}{}{}}
\makeatother
\crefname{originallemma}{original lemma}{original lemmas}
\Crefname{originallemma}{Original lemma}{Original lemmas}
\makeatletter
\newlabel{varieties-lemma-regular-functions-proper-variety}{{04L2}{source}{}{sourcepointer-5}{}}
\newlabel{varieties-lemma-regular-functions-proper-variety@cref}{{[originallemma][0][]04L2}{[1][0][]source}{}{}{}}
\makeatother
\title{The genus of the curve in terms of the numerical type of the special fibre}
\author{The Stacks Project Authors}
\date{stacks-project:models.tex:local-snapshot-20261012; GNU FDL 1.2 or later}
\begin{document}
\maketitle
\noindent\textit{Editorial note.} The counted unit is the Stacks Project Lemma lemma-genus-formula at source lines 5156-5165 together with its proof at source lines 5167-5219, a straight adjacent pairing in the chapter file. The chapter is the upstream file models.tex; the pool directory carries the longer name stacks-models and is not used as the chapter name anywhere in this card.\par\medskip
\noindent\textit{Editorial note.} This card has no context span: the statement and proof are self-contained apart from the five pointers, and the notions they use are recorded in the standard prerequisites.\par\medskip
\noindent\textit{Editorial note.} The retained proof names five results: the situation that the statement refers to, two lemmas of this chapter and one each of the Perfect and Varieties chapters. All five are declared as non-load-bearing navigation pointers with their stable Stacks tags looked up in the shared tag index, and none of their statements is reproduced inside the span. The situation is declared with the kind section, because the assembler's whitelist of kinds has no entry for a situation.\par\medskip
\noindent\textit{Editorial note.} The retained proof span starts at line 5167, the line of the opening marker read from the source, and contains exactly one proof environment.\par\medskip
\noindent\textit{Editorial note.} No citation occurs in the retained spans, so no bibliography entry is supplied.\par\medskip
\noindent\textit{Allowed standard prerequisites (editorial).}\begin{itemize}
\item The regular-model situation, the numerical invariants of the components and the Euler characteristic of a curve. The situation fixes a discrete valuation ring, a smooth projective curve over its fraction field, a regular proper model and the irreducible components of the special fibre with their multiplicities; the component carries a residue field, a genus and a self-intersection number, all recorded by the numerical type. The Euler characteristic of a coherent sheaf on a proper scheme over a field is additive in short exact sequences and, for a family of curves, its value is locally constant in the geometric sense; on a reduced curve it equals the difference between the dimension of the zeroth and first cohomology. The author's own definitions and lemmas of the chapter and the quoted results of the Perfect and Varieties chapters are the allowed inputs.. The counted lemma is the formula relating the genus of the curve to the numerical type of the special fibre, and the chapter's later comparisons of genera of different models are applications of it.
\end{itemize}\medskip
\ifdefined\hypertarget\hypertarget{sourcepointer-1}{}\fi
\noindent\textit{Source reference.} Original section 0C61: focused-seeds/AGP/sources/stacks-snapshot-20261011/tags.txt, source lines 15281--15281. Pointer to the situation that the retained statement names, declared with the kind section because the assembler's whitelist of kinds has no entry for a situation. The number is the stable Stacks tag of the target, looked up in the shared tag index.\par\medskip
\ifdefined\hypertarget\hypertarget{sourcepointer-2}{}\fi
\noindent\textit{Source reference.} Original lemma 0CA2: focused-seeds/AGP/sources/stacks-snapshot-20261011/tags.txt, source lines 15422--15422. Navigation pointer invoked by the retained proof for the way in which the structure sheaf of the special fibre relates to those of its components; the cited statement is an allowed input of the same section and its proof is not reproduced here. The number is the stable Stacks tag of the target, looked up in the shared tag index.\par\medskip
\ifdefined\hypertarget\hypertarget{sourcepointer-3}{}\fi
\noindent\textit{Source reference.} Original lemma 0C66: focused-seeds/AGP/sources/stacks-snapshot-20261011/tags.txt, source lines 15286--15286. Navigation pointer invoked by the retained proof for the numerical properties of the intersection form of a numerical type; the cited statement is an allowed input of the same chapter and its proof is not reproduced here. The number is the stable Stacks tag of the target, looked up in the shared tag index.\par\medskip
\ifdefined\hypertarget\hypertarget{sourcepointer-4}{}\fi
\noindent\textit{Source reference.} Original lemma 0B9T: focused-seeds/AGP/sources/stacks-snapshot-20261011/tags.txt, source lines 14186--14186. Allowed input of the Perfect chapter quoted by the retained proof for the local constancy of the Euler characteristic in a family of curves; the pointer is registered so the reference resolves and the cited statement is not reproduced inside the span. The number is the stable Stacks tag of the target, looked up in the shared tag index.\par\medskip
\ifdefined\hypertarget\hypertarget{sourcepointer-5}{}\fi
\noindent\textit{Source reference.} Original lemma 04L2: focused-seeds/AGP/sources/stacks-snapshot-20261011/tags.txt, source lines 5993--5993. Allowed input of the Varieties chapter quoted by the retained proof for the identification of the residue field of a component with the global sections of its structure sheaf; the pointer is registered so the reference resolves and the cited statement is not reproduced inside the span. The number is the stable Stacks tag of the target, looked up in the shared tag index.\par\medskip
\begin{lemma}
\label{lemma-genus-formula}
In Situation \ref{situation-regular-model} we have
$$
g_C = 1 + \sum\nolimits_{i = 1, \ldots, n}
m_i\left([\kappa_i : k] (g_i - 1) - \frac{1}{2}(C_i \cdot C_i)\right)
$$
where $\kappa_i = H^0(C_i, \mathcal{O}_{C_i})$,
$g_i$ is the genus of $C_i$, and $g_C$ is the genus of $C$.
\end{lemma}

\begin{proof}
Our basic tool will be Derived Categories of Schemes, Lemma
\ref{perfect-lemma-chi-locally-constant-geometric}
which shows that
$$
1 - g_C = \chi(C, \mathcal{O}_C) =
\chi(X_k, \mathcal{O}_{X_k})
$$
Choose a sequence of effective Cartier divisors
$$
X_k = D_m \supset D_{m - 1} \supset \ldots \supset D_1 \supset D_0 = \emptyset
$$
such that $D_{j + 1} = D_j + C_{i_j}$ for each $j$. (It is clear that
we can choose such a sequence by decreasing one nonzero multiplicity
of $D_{j + 1}$ one step at a time.) Applying Lemma \ref{lemma-add-component}
starting with $\chi(\mathcal{O}_{D_0}) = 0$ we get
\begin{align*}
1 - g_C
& =
\chi(X_k, \mathcal{O}_{X_k}) \\
& =
\sum\nolimits_j
\left(-(D_j \cdot C_{i_j}) +  \chi(C_{i_j}, \mathcal{O}_{C_{i_j}})\right) \\
& =
- \sum\nolimits_j
(C_{i_1} + C_{i_2} + \ldots + C_{i_{j - 1}} \cdot C_{i_j}) +
\sum\nolimits_j \chi(C_{i_j}, \mathcal{O}_{C_{i_j}}) \\
& =
-\frac{1}{2}\sum\nolimits_{j \not = j'} (C_{i_{j'}} \cdot C_{i_j}) +
\sum m_i \chi(C_i, \mathcal{O}_{C_i}) \\
& =
\frac{1}{2} \sum m_i(C_i \cdot C_i) + \sum m_i \chi(C_i, \mathcal{O}_{C_i})
\end{align*}
Perhaps the last equality deserves some explanation. Namely, since
$\sum_j C_{i_j} = \sum m_i C_i$ we have
$(\sum_j C_{i_j} \cdot \sum_j C_{i_j}) = 0$ by
Lemma \ref{lemma-properties-form}. Thus we see that
$$
0 = \sum\nolimits_{j \not = j'} (C_{i_{j'}} \cdot C_{i_j}) +
\sum m_i(C_i \cdot C_i)
$$
by splitting this product into ``nondiagonal'' and ``diagonal'' terms.
Note that $\kappa_i$ is a field finite over $k$ by
Varieties, Lemma \ref{varieties-lemma-regular-functions-proper-variety}.
Hence the genus of $C_i$ is defined and we have
$\chi(C_i, \mathcal{O}_{C_i}) = [\kappa_i : k](1 - g_i)$.
Putting everything together and rearranging terms we get
$$
g_C = - \frac{1}{2}\sum m_i(C_i \cdot C_i) +
\sum m_i[\kappa_i : k](g_i - 1) + 1
$$
which is what the lemma says too.
\end{proof}

\end{document}
